% ============================================================
%  CHAPITRE 5 : Benchmark des logiciels de modélisation 3D
% ============================================================
\chapter{Benchmark des logiciels de modélisation 3D}
\label{chap:benchmark_logiciels}

% SECTION 5.1 : INTRODUCTION
\section{Introduction}
\label{sec:introduction_benchmark}

La modélisation 3D des écoulements dans les ouvrages hydrauliques est devenue un outil essentiel pour l'analyse et l'optimisation des performances. De nombreux logiciels CFD sont disponibles, chacun présentant des caractéristiques spécifiques en termes de méthodes numériques, de modèles physiques, de type de maillage et d'accessibilité (open source ou commercial). Ce chapitre propose une analyse comparative des principales solutions logicielles adaptées à la modélisation des écoulements à surface libre, afin de justifier le choix d'OpenFOAM pour la présente étude.

% SECTION 5.2 : PRÉSENTATION DES LOGICIELS
\section{Présentation des logiciels}
\label{sec:presentation_logiciels}

\subsection{TELEMAC-3D}
\label{subsec:telemac}

TELEMAC-3D est un code de calcul open source développé par le Laboratoire National d'Hydraulique et Environnement (LNHE) d'EDF R\&D. Il résout les équations de Navier-Stokes pour les écoulements à surface libre, en utilisant la méthode des \textbf{éléments finis}. Il est particulièrement adapté à l'hydrodynamique côtière, estuarienne et fluviale, ainsi qu'au transport de sédiments et à la qualité de l'eau. Sa flexibilité et sa gratuité en font un outil académique et professionnel de premier plan, mais sa prise en main peut être complexe \cite{galland1991telemac, hervouet2007hydrodynamics}.

\subsection{Flow-3D}
\label{subsec:flow3d}

Flow-3D est un logiciel commercial CFD développé par Flow Science. Il est réputé pour sa capacité à modéliser avec précision les écoulements à surface libre grâce à la méthode \textbf{VOF (Volume of Fluid)}, dont il est l'un des pionniers \cite{hirt1981volume}. Il utilise un maillage structuré cartésien, ce qui simplifie la génération du maillage. Sa méthode \textbf{FAVOR (Fractional Area/Volume Obstacle Representation)} permet de représenter des géométries complexes sur un maillage simple \cite{flow2016user}. Il est très utilisé en hydraulique (déversoirs, ressauts, dissipateurs d'énergie) mais est payant.

\subsection{ANSYS Fluent}
\label{subsec:ansys_fluent}

Ansys Fluent est l'un des logiciels CFD commerciaux les plus complets et les plus utilisés dans le monde, toutes industries confondues (aéronautique, automobile, énergie, hydraulique). Il utilise la méthode des \textbf{volumes finis} sur des maillages non structurés (polyédriques, tétraédriques, hexaédriques). Il dispose d'une très large bibliothèque de modèles physiques (turbulence, multiphasique, combustion, etc.) et d'une interface utilisateur graphique très avancée \cite{ansys2020fluent, matsson2022introduction}. Sa puissance a un coût élevé en termes de licence.

\subsection{XFlow}
\label{subsec:xflow}

XFlow est un logiciel CFD commercial de nouvelle génération développé par Dassault Systèmes. Il se distingue par son approche innovante : il n'utilise pas de maillage traditionnel mais une méthode particulaire basée sur la \textbf{méthode de Lattice Boltzmann} \cite{chen1998lattice, kruger2017lattice}. Cette approche sans maillage est très efficace pour traiter des géométries complexes en mouvement et des écoulements instationnaires. Il est plus récent et son application en hydraulique "classique" est en développement.

\subsection{STAR-CCM+}
\label{subsec:starccm}

STAR-CCM+ est un logiciel commercial CFD développé par Siemens. Il est reconnu pour son interface unifiée et son workflow intégré, de la préparation de la géométrie à l'analyse des résultats. Il utilise également la méthode des volumes finis et a popularisé l'utilisation de maillages polyédriques, qui offrent un bon compromis entre précision et nombre de cellules \cite{cd2019star, peric2005advantage}. Il est très présent dans les secteurs automobile et aéronautique.

\subsection{OpenFOAM}
\label{subsec:openfoam}

OpenFOAM (Open Field Operation and Manipulation) est la principale solution \textbf{open source} en CFD. Développé par OpenCFD Ltd, il s'agit d'une bibliothèque de solveurs et d'outils écrits en C++, permettant de résoudre une vaste gamme de problèmes en mécanique des fluides (et des solides) \cite{weller1998tensorial, jasak2009openfoam}. Il utilise la méthode des volumes finis sur des maillages non structurés. Son principal avantage est sa gratuité, sa transparence (code source accessible) et sa flexibilité : l'utilisateur peut modifier les solveurs existants ou en créer de nouveaux. L'inconvénient est l'absence d'interface graphique native et une courbe d'apprentissage plus raide.

\subsection{SimScale}
\label{subsec:simscale}

SimScale est une plateforme de simulation \textbf{cloud} (SaaS - Software as a Service). Elle propose un accès à des solveurs de CFD (basés sur OpenFOAM) et d'analyse par éléments finis via un navigateur web \cite{simscale2021documentation, bletzinger2017cloud}. Elle fonctionne sur un modèle freemium (compte gratuit avec des limitations) ou par abonnement. Son avantage est l'accessibilité sans nécessiter de matériel puissant, mais la confidentialité des données peut être une préoccupation.

\subsection{Autres logiciels}
\label{subsec:autres_logiciels}

D'autres logiciels complètent le paysage des solutions disponibles :
\begin{itemize}
    \item \textbf{Autodesk CFD} : Solution commerciale d'Autodesk, basée sur la méthode des éléments finis, intégrée dans l'écosystème de conception.
    \item \textbf{ANSYS CFX} : Autre solveur CFD majeur d'Ansys, réputé pour sa robustesse et sa précision, notamment pour les turbomachines.
    \item \textbf{MIKE 3 Wave FM} : Logiciel commercial de DHI, spécialisé dans l'hydrodynamique côtière et maritime (vagues, courants) \cite{dhi2020mike}.
    \item \textbf{COMSOL Multiphysics} : Plateforme commerciale permettant de coupler la CFD avec d'autres physiques (électromagnétisme, thermique, mécanique des structures) \cite{comsol2021reference}.
\end{itemize}

% SECTION 5.3 : COMPARAISON ENTRE LES LOGICIELS
\section{Comparaison entre les logiciels}
\label{sec:comparaison_logiciels}

Le tableau suivant regroupe les principales caractéristiques des logiciels présentés :

\begin{table}[H]
    \centering
    \caption{Comparaison des principaux logiciels de CFD}
    \label{tab:comparaison_logiciels_cfd}
    \resizebox{\textwidth}{!}{%
    \begin{tabular}{|l|c|c|c|c|c|}
        \hline
        \textbf{Logiciel} & \textbf{Interface graphique} & \textbf{Méthode} & \textbf{Maillage 3D} & \textbf{Open source} & \textbf{Coût} \\
        \hline
        TELEMAC-3D & Oui & Éléments finis & Non structuré & \textbf{Oui} & \textbf{Gratuit} \\
        \hline
        Flow-3D & Oui & Volumes finis & Cartésien structuré & Non & Payant \\
        \hline
        ANSYS Fluent & Oui & Volumes finis & Non structuré (polyédrique) & Non & Payant \\
        \hline
        XFlow & Oui & Lattice Boltzmann & Sans maillage & Non & Payant \\
        \hline
        STAR-CCM+ & Oui & Volumes finis & Non structuré (polyédrique) & Non & Payant \\
        \hline
        \textbf{OpenFOAM} & \textbf{Non} (via ParaView) & Volumes finis & Non structuré & \textbf{Oui} & \textbf{Gratuit} \\
        \hline
        SimScale & Oui (cloud) & Volumes finis / EF & Automatique (cloud) & Basé sur OpenFOAM & Freemium \\
        \hline
        ANSYS CFX & Oui & Volumes finis (couplé) & Non structuré & Non & Payant \\
        \hline
    \end{tabular}}
\end{table}

% SECTION 5.4 : LE LOGICIEL SKETCHUP
\section{Le logiciel SketchUp}
\label{sec:sketchup}

Bien que n'étant pas un logiciel de CFD, SketchUp joue un rôle important dans la chaîne de modélisation pour la création de la géométrie. Dans notre démarche, les profils types 2D des ouvrages ont d'abord été réalisés sous \textbf{AutoCAD}, puis importés dans SketchUp afin de servir de base à la modélisation 3D.

\subsection{Présentation du logiciel}
\label{subsec:presentation_sketchup}

SketchUp est un logiciel de modélisation 3D développé par Trimble Inc., reconnu pour la simplicité et l'intuitivité de son interface \cite{trimble2023sketchup, brightman2018sketchup}. Contrairement aux logiciels BIM paramétriques (comme Revit), SketchUp repose sur une modélisation directe à base de faces et d'arêtes, manipulées à l'aide d'outils simples tels que ``Pousser/Tirer'' (Push/Pull), ``Suivez-moi'' ou encore les outils de mise à l'échelle et de rotation. Les éléments peuvent être organisés en \textbf{groupes} et \textbf{composants}, ce qui permet de structurer le modèle et de réutiliser des éléments identiques.

\subsection{Intérêt pour la modélisation des ouvrages hydrauliques}
\label{subsec:interet_sketchup}

Pour la modélisation des ouvrages hydrauliques comme le barrage Ahl Souss, SketchUp présente plusieurs avantages significatifs :

\begin{itemize}
    \item \textbf{Rapidité et simplicité de modélisation} : les outils de modélisation directe (Pousser/Tirer, Suivez-moi) permettent de générer rapidement des géométries complexes (corps du barrage, seuil Creager, coursier, conduites de vidange, murs bajoyers, prises d'eau) à partir de profils 2D \cite{brightman2018sketchup}.
    \item \textbf{Modélisation du terrain naturel} : l'outil \textit{Sandbox} de SketchUp permet de générer une surface topographique (TIN) directement à partir des courbes de niveau, facilitant l'intégration du terrain naturel dans le modèle global.
    \item \textbf{Organisation en groupes et composants} : chaque élément de l'ouvrage peut être isolé dans un groupe ou un composant dédié, ce qui simplifie les modifications et la gestion des différentes parties du modèle.
    \item \textbf{Exportation vers les formats CAO standards} : grâce à l'extension \textbf{SketchUp STL}, les modèles 3D peuvent être exportés au format \textbf{STL (Standard Tessellation Language)} \cite{lhuillier2014sketchupstl}. Ce format, qui décrit une géométrie par un maillage de triangles, est compatible avec la plupart des logiciels de CFD, dont OpenFOAM.
    \item \textbf{Accessibilité} : SketchUp est disponible en version gratuite (SketchUp Free, en ligne) et en version professionnelle à coût modéré, ce qui en fait un outil accessible pour des projets académiques.
\end{itemize}

\subsection{Démarche de modélisation}
\label{subsec:demarche_sketchup}

Le processus de modélisation pour notre étude comprend les étapes suivantes :
\begin{enumerate}
    \item Réalisation du profil type 2D du barrage (cote de fondation, fruit des parements, largeur en crête) et de l'évacuateur de crues sous \textbf{AutoCAD}.
    \item Importation du profil 2D AutoCAD dans un projet SketchUp avec les unités appropriées (mètres).
    \item Modélisation du terrain naturel à l'aide de l'outil \textit{Sandbox} à partir des données topographiques (courbes de niveau).
    \item Création du corps du barrage par extrusion (Pousser/Tirer) du profil type importé depuis AutoCAD.
    \item Modélisation de l'évacuateur de crues (seuil Creager, coursier, cuillère) à l'aide des outils de courbes et de l'outil \textit{Suivez-moi}.
    \item Modélisation de la vidange de fond (conduites, chambre des vannes, bassin de dissipation).
    \item Modélisation des prises d'eau et de la galerie.
    \item Exportation au format STL via l'extension SketchUp STL pour une utilisation dans OpenFOAM \cite{lhuillier2014sketchupstl}.
\end{enumerate}

% SECTION 5.5 : CONCLUSION
\section{Conclusion}
\label{sec:conclusion_benchmark}

Le choix d'OpenFOAM pour cette étude se justifie par plusieurs avantages : sa gratuité, l'accès au code source permettant une transparence totale et une flexibilité maximale, sa capacité de calcul parallèle (HPC) pour traiter des géométries complexes, et sa large communauté d'utilisateurs qui garantit un support et un partage de connaissances. Sa structure modulaire permet d'adapter précisément les modèles physiques aux spécificités des ouvrages du barrage Ahl Souss. Associé à SketchUp pour la modélisation géométrique 3D, cet ensemble d'outils constitue une chaîne de simulation complète, performante et économiquement accessible.

